\section{Detection}\label{sec:detection}

\subsection{Background Subtraction (MOG2 + CLAHE)}
As a baseline, we adopt MOG2 background subtraction \cite{survey}, in which each
pixel's intensity distribution is modelled as a mixture of $K$ Gaussians:
\begin{equation}
  P(x_t) = \sum_{k=1}^{K} w_{k,t}\,\mathcal{N}(x_t \mid \mu_{k,t},\,\sigma_{k,t}^2).
\end{equation}
A pixel is classified as foreground if its Mahalanobis distance to the
best-matching background component exceeds a fixed threshold.
Prior to background modelling, illumination is normalised with Contrast Limited
Adaptive Histogram Equalization (CLAHE) \cite{clahe}, and the resulting foreground
mask is cleaned with morphological opening and closing; connected components smaller
than a minimum area threshold are discarded.
Alternative background subtractors (KNN, LSBP, GSOC) were evaluated but offered no
improvement over MOG2 while incurring higher computational cost.

Despite these refinements, the method exhibits three recurring failure modes: floor
specularities produce shadows whose hue values fall outside the range that MOG2's
built-in shadow detector handles; players who are briefly stationary are
indistinguishable from background; and moving spectators generate persistent false
positives.
Because the optimal morphological kernel sizes and detection thresholds differ
across cameras, the method also lacks cross-camera generalization.

\subsection{Pre-Trained YOLOv11}
To overcome these limitations, we apply a YOLOv11 model \cite{yolo_paper,yolo_github}
pre-trained on COCO, using a confidence threshold of 0.4, an NMS IoU threshold of
0.45, and an inference resolution of 640\,px.
Detection quality improves substantially, with fewer false positives and better
spatial localisation.
However, the pre-trained model has two domain gaps: the ball is absent from COCO's
person-centric classes, and players and referees are both labelled ``person'',
making identity distinction impossible.

\subsection{Fine-Tuned YOLOv26 with Multi-Pass Inference}
To close both gaps, we fine-tune the YOLOv26 backbone on a Sanbapolis-specific
annotated set.
The fine-tuned model uses identity-encoding class labels (one class per player
identity, e.g.\ \textit{Red\_11}) and a dedicated \textit{Ball} class.

Because at 640\,px the convolutional backbone downsamples a 1080p frame by
approximately $1.7\times$, a 15\,px basketball occupies fewer than two feature-map
grid cells, which is insufficient for reliable localisation.
We therefore run two separate inference passes per frame:
\begin{itemize}
  \item \textbf{Player pass}: 1280\,px resolution, confidence threshold 0.3.
    A second run at a lower threshold (0.15) retains a low-confidence band that is
    reserved for the tracker's secondary association pass.
  \item \textbf{Ball pass}: 1600\,px resolution, confidence threshold 0.2.
    The higher resolution increases the spatial footprint of the ball, while the
    lower threshold favours recall for this motion-blurred, occasionally occluded object.
\end{itemize}
The two sets of detections are merged frame by frame into a single output.

\subsubsection{Class-Agnostic NMS}
Because player identity is encoded directly in the class label, the standard
per-class NMS may retain two overlapping boxes that predict different identities for
the same physical player.
We therefore replace it with a class-agnostic variant \cite{nms}: all detections are
pooled regardless of class, sorted by confidence, and greedily suppressed so that
any box whose Intersection over Union with a higher-scoring box exceeds 0.75 is
discarded.
This collapses duplicate boxes for the same player without affecting the ball.
The result of this step is the final detection output fed to the tracker.

\begin{figure}[t]
  \centering
  \includegraphics[width=0.95\linewidth]{../media/fig_detection.png}
  \caption{Sample detection output from fine-tuned YOLOv26 with class-agnostic NMS
           on cam\_13 (three selected frames). Colour-coded bounding boxes indicate
           player identities; the ball is shown separately.}
  \label{fig:detection}
\end{figure}
